% =====================================================================
\section{Exercise 1 -- Bivariate copula simulation, dependence, and IFM/CML}
\label{sec:ex1}
% =====================================================================

We work with six bivariate copulas (\textsc{Clayton} $\theta=10$, \textsc{Survival
Clayton} $\theta=5$, \textsc{Gaussian} $\rho=-0.7$, \textsc{Student-$t$}
$\rho=-0.7,\nu=5$, \textsc{FGM} $\lambda=-0.5$ and \textsc{FGM} $\lambda=0$).
The implementation lives in \file{exercise1/} and is driven by
\file{exercise1/main.py}. Throughout, $V_1,V_2$ denote independent
$\mathcal{U}(0,1)$ draws, while $(U_1,U_2)$
is the resulting copula sample.

% ---------------------------------------------------------------
\subsection{Theoretical building block: the conditional distribution method}
% ---------------------------------------------------------------

Given a bivariate copula $C$, define the partial derivative
\[
C_2(u_1\given u_2)\;=\;\Pr(U_1\le u_1\mid U_2=u_2)\;=\;\frac{\partial}{\partial u_2}\,C(u_1,u_2).
\]
By construction $C_2(\cdot\given u_2)$ is a CDF on $(0,1)$. The
\emph{conditional inversion algorithm} produces a sample from $C$ in two steps:
\[
\boxed{\;U_2\sim\mathcal{U}(0,1),\qquad U_1\;=\;C_2^{-1}(V_1\given U_2),\quad V_1\sim\mathcal{U}(0,1).\;}
\]
This is what every routine in \file{exercise1/copula\_sim.py} implements.

% ---------------------------------------------------------------
\subsection{Closed-form simulators}
% ---------------------------------------------------------------

\paragraph{Clayton ($\theta>0$).}
$C^{Cla}(u_1,u_2)=\bigl(u_1^{-\theta}+u_2^{-\theta}-1\bigr)^{-1/\theta}$.
Differentiating in $u_2$,
\[
C_2(u_1\given u_2)=u_2^{-\theta-1}\bigl(u_1^{-\theta}+u_2^{-\theta}-1\bigr)^{-1-1/\theta}.
\]
Setting $C_2=v_1$ and solving for $u_1$ yields
\begin{equation}\label{eq:clayton-sim}
U_1=\Bigl[\,1+u_2^{-\theta}\bigl(v_1^{-\theta/(\theta+1)}-1\bigr)\Bigr]^{-1/\theta}.
\end{equation}

\paragraph{Survival Clayton.} The survival copula
$\hat C(u_1,u_2)=u_1+u_2-1+C(1-u_1,1-u_2)$ is sampled by drawing
$(\tilde U_1,\tilde U_2)$ from Clayton and setting
$U_i=1-\tilde U_i$, which produces the upper-tail counterpart of the Clayton.

\paragraph{Gaussian ($-1<\rho<1$).}
With $X_i=\Phi^{-1}(U_i)$,
\[
X_1\given X_2=x_2\;\sim\;\mathcal{N}\!\bigl(\rho x_2,\;1-\rho^2\bigr),
\]
so the conditional CDF is $C^G(u_1\given u_2)=
\Phi\!\bigl((\Phi^{-1}u_1-\rho\,\Phi^{-1}u_2)/\sqrt{1-\rho^2}\bigr)$
and the inverse gives
\begin{equation}\label{eq:gauss-sim}
U_1=\Phi\!\bigl(\rho\,\Phi^{-1}(u_2)+\sqrt{1-\rho^2}\,\Phi^{-1}(v_1)\bigr).
\end{equation}

\paragraph{Student $t$ ($\rho,\nu$).} Letting $y_i=t_\nu^{-1}(u_i)$,
\[
\frac{\,Y_1-\rho Y_2\,}{\sqrt{(1-\rho^2)(\nu+Y_2^2)/(\nu+1)}}\;\sim\;t_{\nu+1},
\]
which inverts to
\begin{equation}\label{eq:t-sim}
U_1=t_\nu\!\Bigl(\rho\, Y_2+t_{\nu+1}^{-1}(v_1)\sqrt{(1-\rho^2)(\nu+Y_2^2)/(\nu+1)}\Bigr).
\end{equation}

\paragraph{FGM ($-1\le\lambda\le 1$).}
$C^{FGM}(u_1,u_2)=u_1u_2[1+\lambda(1-u_1)(1-u_2)]$. Setting
$\psi=\lambda(1-2u_2)$, the equation $C_2(u_1\given u_2)=v_1$ becomes
$\psi u_1^2-(1+\psi)u_1+v_1=0$, with the admissible root
\begin{equation}\label{eq:fgm-sim}
U_1=\frac{(1+\psi)-\sqrt{(1+\psi)^2-4\psi v_1}}{2\psi}\,.
\end{equation}
At $\lambda=0$ we have independence and $U_1=v_1$ is recovered as the limit.

\paragraph{Mixture (used in 1.d).} A copula mixture
$C=\pi C_A+(1-\pi)C_B$ is sampled by drawing a Bernoulli($\pi$) indicator
and selecting $C_A$ or $C_B$.

% ---------------------------------------------------------------
\subsection{1.a -- Scatter plots and dependence measures}
% ---------------------------------------------------------------

For each copula we draw $N=10^5$ pairs and compute the three
dependence measures: Pearson's $\rho$, Spearman's $\rho_S$ and Kendall's $\tau$.
Pearson is sensitive to the marginals and an average, while Spearman and Kendall 
capture better tail dependence.
Numerical results (Table~\ref{tab:1a}) and scatter plots
(Figure~\ref{fig:1a}) show:
\begin{itemize}\setlength\itemsep{0pt}
\item Clayton has lower-tail dependence.
\item Survival Clayton has upper-tail dependence.
\item Gaussian and Student-$t$ at $\rho=-0.7$ produce the elliptical dependence;
      $t$ adds heavier joint extremes.
\item FGM is essentially close to independence: its dependence parameter
      is bounded ($|\tau|\le 2/9$).
\end{itemize}

\begin{table}[H]\centering\small
\caption{Dependence measures, $N=10^5$ draws.}\label{tab:1a}
\begin{tabular}{lrrr}
\toprule
Copula & Pearson & Spearman & Kendall\\\midrule
Clayton $\theta=10$            & $0.9584$ & $0.9583$ & $0.8335$\\
Survival Clayton $\theta=5$    & $0.8844$ & $0.8847$ & $0.7141$\\
Gaussian $\rho=-0.7$           & $-0.6840$ & $-0.6839$ & $-0.4947$\\
Student-$t$ $\rho=-0.7,\nu=5$  & $-0.6723$ & $-0.6722$ & $-0.4948$\\
FGM $\lambda=-0.5$             & $-0.1706$ & $-0.1706$ & $-0.1139$\\
FGM $\lambda=0$                & $-0.0044$ & $-0.0044$ & $-0.0029$\\
\bottomrule
\end{tabular}
\end{table}

\begin{figure}[H]\centering
\includegraphics[width=0.95\linewidth]{fig_1a_scatter.png}
\caption{Dependent $\mathcal{U}(0,1)$ scatter plots, $N=10^5$.}\label{fig:1a}
\end{figure}

\paragraph{Kendall's $\tau$:} For Clayton,
$\tau=\theta/(\theta+2)=10/12=0.8333$, matching the simulated $0.8335$. For
Gaussian/$t$, $\tau=(2/\pi)\arcsin\rho=-0.4936$, matching $-0.4947$ and $-0.4948$.

% ---------------------------------------------------------------
\subsection{1.b -- Portfolio statistics, VaR, and correlation across copulas}
% ---------------------------------------------------------------

We build the portfolio $r_p=0.4\,r_1+0.6\,r_2$ where
$r_1=\Phi^{-1}(U_1)\sim\mathcal{N}(0,1)$ and
$r_2=t_6^{-1}(U_2)\sim t_6$. The marginals enter only through their quantile
functions; the copula governs the joint behaviour. Statistics are reported
in Table~\ref{tab:1b}, and a direct comparison of histograms vs.\ a Normal
overlay (same $\mu,\sigma$) is shown in Figure~\ref{fig:1b}.

\begin{table}[H]\centering\small
\caption{Portfolio statistics for $r_p=0.4r_1+0.6r_2$, $r_1\!\sim\!\mathcal{N}(0,1)$,
$r_2\!\sim\!t_6$.}\label{tab:1b}
\begin{tabular}{lrrrrrr}
\toprule
Copula & Mean & Std & Skew & Kurt & VaR$_{1\%}$ & VaR$_{5\%}$\\\midrule
Clayton           & $0.0007$ & $1.1116$ & $-0.1683$ & $4.1135$ & $-2.7795$ & $-1.8168$\\
Surv. Clayton     & $0.0012$ & $1.0915$ & $\phantom{-}0.2290$ & $4.0641$ & $-2.4421$ & $-1.6886$\\
Gaussian          & $0.0007$ & $0.5394$ & $-0.0197$ & $6.6978$ & $-1.3527$ & $-0.8432$\\
Student-$t$       & $0.0008$ & $0.5399$ & $\phantom{-}0.0029$ & $6.8199$ & $-1.3815$ & $-0.8472$\\
FGM $\lambda=-0.5$ & $0.0010$ & $0.7786$ & $-0.0039$ & $5.1494$ & $-1.9349$ & $-1.2390$\\
FGM $\lambda=0$    & $0.0011$ & $0.8342$ & $-0.0023$ & $4.7379$ & $-2.0458$ & $-1.3340$\\
\bottomrule
\end{tabular}
\end{table}

\begin{figure}[H]\centering
\includegraphics[width=0.95\linewidth]{fig_1b_histograms.png}
\caption{Portfolio return histograms with Normal($\mu,\sigma^2$) overlay.}\label{fig:1b}
\end{figure}

\paragraph{Reading the table.}
The variance of $r_p$ is
\[
\Var(r_p)=w_1^2\Var(r_1)+w_2^2\Var(r_2)+2w_1w_2\Cov(r_1,r_2),
\]
so a positive copula increases $\Var(r_p)$ relative to a negatively dependent
one (compare Clayton vs.\ Gaussian). VaR at level $\alpha$ is the empirical
$\alpha$-quantile $\widehat{\mathrm{VaR}}_\alpha=\hat F_{r_p}^{-1}(\alpha)$
Negative dependence (Gaussian $\rho=-0.7$) drastically reduces $\sigma_p$ and
hence the magnitude of VaR, as expected from diversification.

\newpage
% ---------------------------------------------------------------
\subsection{1.c -- Estimating a Gaussian copula by IFM with EDF margins}
% ---------------------------------------------------------------

Given an i.i.d. sample $\{(r_{1t},r_{2t})\}_{t=1}^T$ the Inference-Functions-for-Margins
(IFM) estimator proceeds in two steps:
\begin{enumerate}\setlength\itemsep{0pt}
\item Estimate each marginal $\hat F_i$ separately;
\item Estimate the copula parameter by maximising the copula log-likelihood
      evaluated at $\hat u_{it}=\hat F_i(r_{it})$.
\end{enumerate}
When $\hat F_i$ is the empirical CDF $\hat F_i(x)=\frac{1}{T+1}\sum_t \indic\{r_{it}\le x\}$
(``EDF'' margin) we obtain the \emph{Canonical Maximum Likelihood} (CML) variant.
The shift from $T$ to $T+1$ avoids values that would blow up
$\Phi^{-1}(\hat u)$.

\paragraph{Gaussian copula log-likelihood (slides eq.\,31).}
Let $\eta_t=(\Phi^{-1}\hat u_{1t},\Phi^{-1}\hat u_{2t})^\top$ and
$\Psi=\bigl(\begin{smallmatrix}1&\rho\\\rho&1\end{smallmatrix}\bigr)$. Then
\begin{align}
\ln c^G(\hat u_{1t},\hat u_{2t};\rho)
&=-\tfrac12\ln(1-\rho^2)
 -\tfrac12\,\eta_t^\top(\Psi^{-1}-I)\eta_t\notag\\
&=-\tfrac12\ln(1-\rho^2)-\frac{\rho^2(\eta_{1t}^2+\eta_{2t}^2)-2\rho\eta_{1t}\eta_{2t}}{2(1-\rho^2)}.
\label{eq:gauss-cop-ll}
\end{align}
We maximise $\sum_t \ln c^G$ in $\rho$ by Brent's method on $(-0.999,0.999)$
(\code{exercise1.estimation}).

\paragraph{Result.} On the Gaussian-copula sample with true $\rho=-0.7$ we obtain
$\hat\rho=-0.702095$, $\ln L=33{,}949.36$, with successful convergence. The
estimator recovers the truth to three decimals on $10^5$ observations.

% ---------------------------------------------------------------
\subsection{1.d -- Mixture DGP and t-copula IFM estimation}
% ---------------------------------------------------------------

We now generate $N=10^5$ pairs from the mixture
\[
C=0.6\;C^{Cla}_{\theta=10}+0.4\;\hat C^{Cla}_{\theta=5},
\]
which produces strong dependence in both tails (Figure~\ref{fig:1d}). The
sample dependence measures are Pearson $0.9283$, Spearman $0.9282$, Kendall $0.7750$.

\begin{figure}[H]\centering
\includegraphics[width=0.65\linewidth]{fig_1d_mixture.png}
\caption{Mixture $0.6\,C^{Cla}_{10}+0.4\,\hat C^{Cla}_{5}$ scatter ($N=10^5$).}
\label{fig:1d}
\end{figure}

\paragraph{Student-$t$ copula log-likelihood (slides eq.\,38).}
Let $y_{it}=t_\nu^{-1}(\hat u_{it})$ and
$Q_t=y_{1t}^2-2\rho y_{1t}y_{2t}+y_{2t}^2$. Then
\begin{align}
\ln c^t(\hat u_{1t},\hat u_{2t};\rho,\nu)
&=\ln\Gamma(\tfrac{\nu+2}{2})+\ln\Gamma(\tfrac{\nu}{2})
   -2\ln\Gamma(\tfrac{\nu+1}{2})-\tfrac12\ln(1-\rho^2)\notag\\
&\quad-\tfrac{\nu+2}{2}\,\ln\!\Bigl(1+\tfrac{Q_t}{\nu(1-\rho^2)}\Bigr)
+\tfrac{\nu+1}{2}\sum_{i=1}^{2}\ln\!\Bigl(1+\tfrac{y_{it}^2}{\nu}\Bigr).
\label{eq:t-cop-ll}
\end{align}
\code{exercise1.estimation.estimate\_t\_copula} maximises this by Nelder--Mead.

\paragraph{Result.} The fitted t-copula returns $\hat\rho=0.929316$, $\hat\nu=2.1698$,
$\ln L=96{,}886.07$. The very small $\hat\nu$ is the t-copula's only way to
mimic the heavy joint tails of the Clayton mixture --- a Gaussian fit on the
same data gives $\hat\rho=0.9057$ and $\ln L=85{,}791$.
This is a textbook example of how IFM correctly
identifies tail dependence by inflating the t-copula degrees of freedom downward.
